
\documentclass[pdflatex,sn-mathphys-num]{sn-jnl}

\DeclareTextCommandDefault{\k}[1]{{\fontencoding{T1}\selectfont\k{#1}}}
\usepackage{graphicx}
\usepackage{amsmath,amssymb,amsfonts}
\usepackage{amsthm}
\usepackage{xcolor}
\usepackage{textcomp}
\usepackage{manyfoot}
\usepackage{booktabs}
\usepackage{algorithm}
\usepackage{algpseudocode}
\usepackage{placeins}

\theoremstyle{thmstyleone}
\newtheorem{theorem}{Theorem}
\newtheorem{lemma}[theorem]{Lemma}
\newtheorem{proposition}[theorem]{Proposition}
\newtheorem{corollary}[theorem]{Corollary}
\theoremstyle{thmstyletwo}
\newtheorem{example}{Example}
\newtheorem{remark}{Remark}
\theoremstyle{thmstylethree}
\newtheorem{definition}{Definition}

\newcommand{\N}{\mathbb{N}}
\newcommand{\Z}{\mathbb{Z}}
\newcommand{\cE}{\mathcal{E}}
\newcommand{\cB}{\mathcal{B}}
\newcommand{\e}{\mathbf{e}}
\newcommand{\out}{\mathrm{out}}
\newcommand{\inn}{\mathrm{in}}
\newcommand{\IN}{\mathrm{IN}}
\newcommand{\OUT}{\mathrm{OUT}}
\newcommand{\FIX}{\textup{\textsc{Fix}}}
\newcommand{\FREE}{\textup{\textsc{Free}}}
\newcommand{\SFIX}{\textup{\textsc{Strict-Fix}}}
\newcommand{\SFREE}{\textup{\textsc{Strict-Free}}}
\newcommand{\abs}[1]{\lvert #1\rvert}
\newcommand{\floor}[1]{\lfloor #1\rfloor}
\newcommand{\YES}{\texttt{YES}}
\newcommand{\NO}{\texttt{NO}}

\raggedbottom

\graphicspath{{figures/}}
\makeatletter
\def\input@path{{sections/}}
\makeatother

\begin{document}

\title[Exact edge cover by temporal walks]{Exact edge cover by strict and non-strict temporal walks}

\author[1]{\fnm{Duc Thien} \sur{Quach}}\email{23020161@vnu.edu.vn}

\author*[1]{\fnm{Manh Cuong} \sur{Tran}}\email{tranmanhcuong@vnu.edu.vn}

\affil[1]{\orgdiv{Faculty of Information Technology}, \orgname{VNU University of Engineering and Technology}, \orgaddress{\city{Hanoi}, \country{Vietnam}}}

\abstract{Exact edge cover asks whether the edges of a temporal graph with $n$ vertices and $m$ edges can be partitioned into time-respecting walks with a prescribed number of walks starting at each vertex and, in the fixed-end variant, ending at each vertex. For non-strict walks, we prove that the problem is polynomially solvable by a new approach based on M-convexity, which also closes a gap in the proof sketch of an earlier static reduction. An orientation of the edges is realizable by walks exactly when walk counters stay non-negative and every snapshot component contains a walk, and only components with all degrees and counters even can violate the latter condition. Once a dead state is removed, the states of each such component form an M-convex set. This yields an augmenting path algorithm that returns walks or an infeasibility certificate, using at most $m/2$ augmentations and $O(n+nm^2)$ arithmetic operations with path-reversing re-orientation. The implementation was validated on small instances. For strict walks, we prove that the problem is NP-complete, already with six time labels for prescribed ends and two for free ends.}

\keywords{Temporal graph, Exact edge cover, Temporal walk, M-convex set, Submodular flow, NP-completeness}

\pacs[MSC Classification]{05C85, 90C27, 90C35, 68Q17}

\maketitle

\input{1-introduction}
\input{2-formulation}
\input{3-orientations}
\input{4-algorithm}
\input{5-strict}
\input{6-experiments}
\input{7-conclusion}
\input{declarations}

\bibliography{references}

\end{document}
